\section{Proofs and Supporting Results}\label{app:proofs}

Throughout, $E(S)=c(S)+n-2r(S)$ and $E_w(S)=\cfp[c(S)-r(S)]+\cfn[n-r(S)]$, with $c$ and $r$ computed from one fixed choice of maximum matching.

\begin{lemma}[Matching size]\label{lem:rank}
Fix the targets, a ground set of candidate records and the compatibility graph between them. Then $S\mapsto r(S)$ is the rank function of a transversal matroid \citep{edmonds1965transversals}: it never decreases when records are added, increases by 0 or 1 when one record is added, and is submodular (adding a record helps a smaller set at least as much as a larger one).
\end{lemma}

\paragraph{Proof of Lemma~\ref{lem:rank}.} Consider the bipartite graph between records and targets with edges given by $\rho_\tau$. The subsets of records that can be matched simultaneously form the independent sets of a transversal matroid \citep{edmonds1965transversals}, and the maximum matching size of $S$ is the rank of $S$ in that matroid. Matroid rank functions are monotone, satisfy $r(S\cup\{z\})-r(S)\in\{0,1\}$, and are submodular. \qed

\paragraph{Proposition~\ref{prop:identity}.} Substituting the definition, $E_w(A)-E_w(S)=\cfp[(c(A)-c(S))-(r(A)-r(S))]-\cfn[r(A)-r(S)]=\cfp\Delta c-(\cfp+\cfn)\Delta r$. Without an ignore region every record is scored, so $\Delta c=|A|-|S|$. \qed

\paragraph{Proposition~\ref{thm:intersect}.} Deleting $A\setminus S$ lowers $E_w$ iff $E_w(A)-E_w(S)>0$, which by Proposition~\ref{prop:identity} is $\cfp\Delta c>(\cfp+\cfn)\Delta r$. If $\Delta c>0$, divide by $(\cfp+\cfn)\Delta c$. If $\Delta c=0$, deletion lowers $E_w$ iff $\Delta r<0$, which cannot happen by monotonicity (Lemma~\ref{lem:rank}), so deletion of no scored record never helps. Addition of $B'$ to $A$ is the same identity with the roles of $A\cup B'$ and $A$: it lowers $E_w$ iff $(\cfp+\cfn)[r(A\cup B')-r(A)]>\cfp[c(A\cup B')-c(A)]$. For the population rule, the expected change is $\cfp\mathbb{E}[\Delta c]-(\cfp+\cfn)\mathbb{E}[\Delta r]$, whose sign is that of $\cfp/(\cfp+\cfn)-p_1$ when $\mathbb{E}[\Delta c]>0$. \qed

\begin{corollary}[F1 under deletion]\label{cor:f1}
If $n+c(A)>0$ and $n+c(S)>0$, then
\[
\mathrm{F1}(S)-\mathrm{F1}(A)=\frac{2\{r(A)\Delta c-[n+c(A)]\Delta r\}}{[n+c(A)][n+c(S)]},
\]
so for $\Delta c>0$ deletion raises F1 on that image iff $\Delta r/\Delta c<\mathrm{F1}(A)/2$. If $n=0$, $c(A)>0$ and every scored record is deleted, F1 rises from 0 to 1.
\end{corollary}
In 22 of our 54 base/partner pairs, the pooled marginal matching rate lies between half the base call's mean F1 and $1/2$; since the condition is per image, this count indicates a direction only. A mean F1 over images changes by the mean of the per-image differences, not by comparing a pooled rate with an average threshold.

\paragraph{Proof of Corollary~\ref{cor:f1}.} Write $r=r(A)$, $c=c(A)$, so that $n+c(S)=n+c-\Delta c$. Then $\mathrm{F1}(S)-\mathrm{F1}(A)=\frac{2(r-\Delta r)}{n+c-\Delta c}-\frac{2r}{n+c}=\frac{2[(r-\Delta r)(n+c)-r(n+c-\Delta c)]}{(n+c)(n+c-\Delta c)}=\frac{2[r\,\Delta c-(n+c)\Delta r]}{(n+c)(n+c-\Delta c)}$. Both denominators are positive, so for $\Delta c>0$ the sign is that of $r/(n+c)-\Delta r/\Delta c=\mathrm{F1}(A)/2-\Delta r/\Delta c$. If $n=0$, then $r=0$ and $\mathrm{F1}(A)=0$ whenever $c(A)>0$; deleting every scored record gives $n=c=0$ and $\mathrm{F1}=1$ by convention. \qed

\paragraph{Proposition~\ref{prop:geometry}.} The decomposition is a telescoping sum. Applying $E=c+n-2r$ to $S_0$ and $S_f$ gives the geometry term $[c(S_0)-c(S_f)]-2[r(S_0)-r(S_f)]$. Convexity of the squared distance gives $\|\bar z-y\|^2\le\frac1m\sum_i\|z_i-y\|^2$ for any target $y$, so averaging does not increase mean squared error to any fixed target. A threshold matching is not monotone in squared error, however. If one member is within tolerance of $y$ and the other is far from every target, their mean can leave the tolerance ball of $y$ and a match is lost; conversely, two members just outside the tolerance on opposite sides of $y$ have a mean inside it, and a match is gained. Hence the geometry term can have either sign. \qed

\paragraph{Proposition~\ref{prop:readout}.} Both errors are evaluated against the same $n$ for the same response, so their difference depends only on which field is read. By the triangle inequality, $\|d-\mathrm{lc}(Z)\|_1\le\|d-n\|_1+\|n-\mathrm{lc}(Z)\|_1\le2\max\{\|d-n\|_1,\|\mathrm{lc}(Z)-n\|_1\}$. For the second inequality, the reverse triangle inequality gives $\big||d_c-n_c|-|\mathrm{lc}(Z)_c-n_c|\big|\le|d_c-\mathrm{lc}(Z)_c|$ for every class $c$; summing over classes and using $|\sum_c x_c|\le\sum_c|x_c|$ gives the bound. \qed

\paragraph{Diminishing coverage.} Let $U_K=Z_1\cup\dots\cup Z_K$ and $\delta_K=\mathbb{E}[r(U_K)-r(U_{K-1})]$. By submodularity, $r(U_K\cup Z_{K+1})-r(U_K)\le r(U_{K-1}\cup Z_{K+1})-r(U_{K-1})$. Taking expectations and using exchangeability of $Z_K$ and $Z_{K+1}$ given the image, the right side has expectation $\delta_K$, so $\delta_{K+1}\le\delta_K$. \qed

\paragraph{Proposition~\ref{prop:partner}.} Under deletion-only confirmation, the records of $A$ not confirmed by $B$ are deleted. In expectation, $n_+(1-h_+^B)$ true and $n_-(1-h_-^B)$ false records are deleted. With no takeover, $\Delta r$ is the number of deleted true records, and with all false records scored, $\Delta c$ is the number of deleted records. By Proposition~\ref{prop:identity}, $\mathbb{E}[E(A)-E(A\cap B)]=n_+(1-h_+^B)+n_-(1-h_-^B)-2n_+(1-h_+^B)=n_-(1-h_-^B)-n_+(1-h_+^B)$. The comparison follows from $G(B)-G(A')=n_-(h_-^{A'}-h_-^{B})-n_+(h_+^{A'}-h_+^{B})$. For the examples, $10(0.15)-100(0.05)=-3.5$ and $10(1)-100(0.5)=-40$. \qed

\paragraph{Proposition~\ref{prop:alarm}.} Let $M_A$ and $M_B$ be maximum matchings of $A$ and $B$ to $Y$ at tolerance $\theta$. At least $r(A)+r(B)-n$ targets are matched in both. For each such target $y$, its partners $a\in A$ and $b\in B$ satisfy $\|a-b\|\le\|a-y\|+\|y-b\|\le2\theta$, and distinct targets give disjoint pairs. Hence $m_{2\theta}(A,B)\ge r(A)+r(B)-n$. With no ignore region, $E(A)+E(B)=|A|+|B|+2n-2r(A)-2r(B)\ge|A|+|B|-2m_{2\theta}(A,B)=U_{2\theta}(A,B)$. The bound on the maximum follows. For boxes, if $\mathrm{IoU}(a,y)\ge t$ and $\mathrm{IoU}(b,y)\ge t$, then $1-\mathrm{IoU}$ is a metric (the Jaccard distance), so $\mathrm{IoU}(a,b)\ge2t-1$, and the same argument applies. At $t=0.5$ the overlap threshold becomes 0, so the bound for boxes keeps differences in the number of records per class but loses all geometric information. \qed

\paragraph{Proposition~\ref{prop:order}.} Every rule $g(\Phi)$ takes values in $\SA$, so pointwise $\ell(g(\Phi),Y)\ge\min_{S\in\SA}\ell(S,Y)$; taking expectations and the infimum over $g$ gives $R_{\mathrm{or}}\le R^*_\Phi$. A fitted rule is one such $g$, so $R^*_\Phi\le R_{\hat g}$. The decomposition is an identity. \qed

\begin{proposition}[Constrained oracle]\label{prop:oracle}
For $\SA$ with at most one representative per cluster, including the option of choosing none, $\min_{S\in\SA}E(S)=n-|M|$, where $M$ is a maximum matching in the bipartite graph between clusters and targets that has an edge whenever some $a\in A(Q)$ is compatible with the target. It is attained by selecting one compatible representative for each matched cluster and nothing for the others.
\end{proposition}
Unrestricted matching over all records can exceed $|M|$ when two members of one cluster match two targets; that value is not attainable within $\SA$ (Table~\ref{tab:select} reports both).

\paragraph{Proof of Proposition~\ref{prop:oracle}.} Any $S\in\SA$ contains at most one record per cluster, so a maximum matching of $S$ to targets induces a cluster-target matching of size $r(S)$ in the stated graph; hence $r(S)\le|M|$. Records unmatched in the chosen maximum matching and lying outside $\Omega$ add to $c$, so $c(S)\ge r(S)$ with or without ignore regions, and $E(S)=c(S)+n-2r(S)\ge n-r(S)\ge n-|M|$. The construction that selects, for each cluster matched in $M$, a representative compatible with its target, and nothing else, has $r=c=|M|$ and attains the bound. \qed

\paragraph{Proposition~\ref{prop:three}.} The output is a representative. If no candidate hits, the output misses (event of $F$). Otherwise, if no representative hits, the output misses (event of $L$). Otherwise some representative hits, and the output misses exactly in the event of $W$. The three events are disjoint and their union is the event that the output misses.

\paragraph{Proposition~\ref{prop:coverage}.} Every candidate of the enlarged menu misses iff every earlier candidate misses and the added one misses, so the probability is $F\,\Pr(\text{added misses}\mid\text{all earlier miss})=F(1-\phi)$ when $F>0$, and $0$ when $F=0$. Under conditional independence given $q$, $\Pr(\text{all miss}\mid q)=\prod_{j=1}^{k}(1-q_j)$; taking expectations gives $F$.

\paragraph{Proposition~\ref{prop:switch}.} The change in the miss indicator is $\mathbb{1}_T\,(\mathbb{1}[b\text{ misses}]-\mathbb{1}[g\text{ misses}])$, which is $+1$ when $b$ misses and $g$ hits, $-1$ when $b$ hits and $g$ misses, and $0$ otherwise. Its expectation is $\Pr(T)[\Pr(b\text{ misses},g\text{ hits}\mid T)-\Pr(b\text{ hits},g\text{ misses}\mid T)]=\Pr(T)[e_Tr_{\mathrm{fix}}-(1-e_T)r_{\mathrm{brk}}]$. If $\mathbb{1}_T$ and $\mathbb{1}[b\text{ misses}]$ are independent given $U$ and $T\subseteq U$, then $\Pr(b\text{ misses}\mid T)=\Pr(b\text{ misses}\mid U,T)=\Pr(b\text{ misses}\mid U)$. The sample size for choosing a scorer follows from Hoeffding's inequality for each scorer's empirical miss rate with a union bound over the $N_{\mathrm{s}}$ candidate scorers.

\paragraph{Proposition~\ref{prop:nocert}.} For any acceptance rule $a(H)\in[0,1]$, the acceptance probability in world $j$ is $\int a\,dP_j$. If $P_0=P_1$ the two are equal, so acceptance at least $1-\beta$ in world 0 forces acceptance at least $1-\beta>\alpha$ in world 1 when $\alpha+\beta<1$. In general, $|\int a\,dP_0-\int a\,dP_1|\le\mathrm{TV}(P_0,P_1)$ for $0\le a\le1$. For the hidden-target construction, couple the two worlds so that the output histories coincide unless the added target is reported in at least one of the $K$ calls, an event of probability at most $K\vartheta$ by the union bound, since the per-call bound $\vartheta$ holds given the earlier history; non-interference means that, conditional on the target not being reported, the rest of the output has the same law in both worlds. The coupling inequality gives $\mathrm{TV}(P_0,P_1)\le K\vartheta$. \qed

\paragraph{Corollary~\ref{cor:nonid}.} By assumption there is an admissible world in which the report over-counts, $|S|>n_0$. Add $m>|S|-n_0$ targets that no interface reports and that do not change the rest of the output ($\vartheta=0$); by closure the result is admissible. The output law is unchanged and $|S|<n_0+m$, so the sign of the count gap differs between two worlds with the same output law. \qed

\subsection{Explanatory parametrization and views}\label{app:param}

To describe how propensities vary with the model and the view when the rest of the interface is held fixed, that is, for calls that share the prompt, schema, setting and cap, we use the explanatory parametrization
\[
\mathrm{logit}\,q_{o,\eta}=a_o+b_{o,m(\eta)}+\gamma_{v(\eta)}(o),\]
\[\textstyle\sum_m b_{o,m}=0,\qquad \gamma_{v_0}(o)=0,
\]
where $v_0$ is the whole image, $a_o$ is common to all models (how hard the object is), $b_{o,m}$ is specific to model $m$, and $\gamma_v(o)$ is the change from the whole image to view $v$, which depends on the apparent size of $o$ in $v$. All three terms may differ between prompts, schemas, settings and caps; a comparison that changes one of these compares two such parametrizations, not two terms of one. The form is related to item-response and labeler-expertise models, in which an item's difficulty and a rater's ability combine on the logit scale \citep{whitehill2009glad} (there as a product, here additively, as in Rasch-type models); here the items are objects and the raters are models and views. The two constraints are conventions that make the parts well defined: centering separates the shared part from the model-specific one, and the reference view separates both from the view effect. We do not estimate the parts; the parametrization motivates the predictions of Section~\ref{sec:theory:behaviour}, each of which is tested directly.

A view can raise how often objects are reported while the misses of the two calls stay as correlated, or become more so. For a true object drawn at random from the image, let $A$ and $A'$ indicate that two calls with the same interface report it and $B$ that a call on another view does, with means $\mu_A=\mu_{A'}$ and $\mu_B$, the calls being conditionally independent given the propensities. Then
\[
\begin{aligned}
&P(A=0,B=1)-P(A=0,A'=1)\\
&\quad=(1-\mu_A)(\mu_B-\mu_A)\\
&\qquad+\mathrm{Cov}(A,A')-\mathrm{Cov}(A,B):
\end{aligned}
\]
the first term is the gain in how often objects are reported, the second the gain in complementarity. For example, with two equally likely kinds of object reported with probability 0.2 and 0.4 on the whole image, a view that adds $\log 2$ to every logit gives $1/3$ and $4/7$: the mean rises from 0.30 to 0.45 and the covariance with the first call rises from 0.0100 to 0.0119, yet the probability that the first call misses an object and the second reports it is 0.305 for the new view against 0.200 for a repeated call: the gain in how often objects are reported (0.107) more than offsets a small loss of complementarity (0.002).
